\documentclass[10pt,a4paper]{amsart}
\usepackage[margin=19mm]{geometry}
\usepackage{amsmath,amssymb,mathtools}
\usepackage[scr=rsfs,cal=euler]{mathalfa}
\usepackage{xfrac}
\usepackage[unicode,hidelinks,bookmarks=false]{hyperref}
\newcommand{\ie}{i.e.\ }
\newcommand{\cf}{cf.\ }
\newcommand{\resp}{respectively\ }
\newcommand{\Subsection}{\subsection}
\providecommand{\nobreakdash}{\nobreak\hbox{-}}
\setlength{\emergencystretch}{2em}
\multlinegap0pt
\usepackage{amssymb}
\newtheorem{lemma}{Lemma}
\title{Nearly invariant subspaces in the real Hardy space}
\author{Arshad Khan, Sneh Lata,, Dinesh Singh}
\date{Source: arXiv:2604.10240v1 (CC BY 4.0)}
\begin{document}
\maketitle

\noindent\emph{Source and licence.} Nearly invariant subspaces in the real Hardy space. Version arXiv:2604.10240v1 (CC BY 4.0), distributed by arXiv under the Creative Commons Attribution 4.0 International licence, http://creativecommons.org/licenses/by/4.0/, as recorded in the arXiv licence metadata for this exact version and shown on the abstract page https://arxiv.org/abs/2604.10240v1. Full licence notice: \texttt{licenses/arxiv-2604.10240-CC-BY-4.0.txt}.\par\smallskip

\noindent\textit{Editorial scope.} Counted unit: the article's lemma Lem3 (source0.tex lines 219--222). The article's complete original proof of it is delivered with it (source0.tex lines 223--255, one \texttt{proof} environment of 33 lines). No mathematics is written, repaired or re-derived: the statement and the proof are the article's own lines, in the article's own notation, and no retained line is substituted or re-flowed.  Retained uncounted as the source-local context the proof needs: lines 185--187.  Nothing was omitted from any retained span, so the package carries no omission record.  No retained line contains a citation, so the wrapper carries no bibliography and no bibliography entry.  No retained line contains a figure environment, an image-inclusion command or drawing code, so no image is delivered and the package declares no asset dependency.  Editorial formatting only, in addition: an AMS \texttt{amsart} preamble that restates the article's own declarations and macros used by the retained lines, namely the environments \texttt{lemma} on separate counters, together with the standard packages those lines need.  The retained environments print in this document as Lemma 1 in order of appearance, whereas the article prints the same items with the numbering of its own full document.  The complete native source of the article as distributed in the arXiv e-print for this exact version is bundled unchanged as \texttt{sources/198253/source0.tex}.
\begin{lemma}\label{Lem1}
Suppose $f$ and $h$ are two functions in $H^2(\mathbb{D})$. Then $$\langle f, h \rangle = \langle \widehat{h}, \widehat{f} \rangle.$$   
\end{lemma}

\begin{lemma}\label{Lem3}
Suppose $\mathcal{N}$ is a non-zero subspace of $H^2(\mathbb{D})$. Then $$H^2_{\mathbb{R}}(\mathbb{D}) \ominus \phi(\mathcal{N})\subset\phi(H^2(\mathbb{D})\ominus \mathcal{N}).$$
Moreover, if $\widehat{\mathcal{N}}\subset \mathcal{N}$, then equality holds.
\end{lemma}

\begin{proof}
Suppose $f \in H^2_\mathbb{R}(\mathbb{D}) \ominus \phi(\mathcal{N})$. Then $$\langle f, \phi(h) \rangle =0  \text{ 
 } \text{for any } h \in \mathcal{N}.$$ 
This implies that for any $h \in \mathcal{N}$,
\begin{eqnarray*}
\langle f, h \rangle &=& -\langle f, \widehat{h} \rangle\\
&=& -\langle h, \widehat{f} \rangle \ \ \  \text{(by Lemma \ref{Lem1})}\\
&=& -\langle h, f \rangle  \ \ \  (\text{since} \ f =\widehat{f}).  
\end{eqnarray*}
Now, consider the following orthogonal decomposition
$$H^2(\mathbb{D})= \mathcal{N} \oplus (H^2(\mathbb{D})\ominus \mathcal{N}).$$
Then, $f$ can be expressed as $f=f_1 + f_2$, where $f_1 \in \mathcal{N}$ and $f_2 \in H^2(\mathbb{D})\ominus \mathcal{N}$. Then
\begin{eqnarray*}
\langle f_1, f_1 \rangle&=&\langle f_1 + f_2, f_1 \rangle\\
&=&\langle f, f_1 \rangle\\
&=&-\langle f_1, f \rangle\\
&=&-\langle f_1, f_1 \rangle ,   
\end{eqnarray*}
which shows that $\|f_1\|=0$; that is, $f_1=0$. Therefore $$f=\phi(f)=\phi(f_2),$$
which belongs to $\phi(H^2(\mathbb{D})\ominus \mathcal{N})$.

 
 Further, if $\widehat{\mathcal{N}}\subset \mathcal{N}$. Choose $f \in \phi(H^2(\mathbb{D})\ominus \mathcal{N})$. Then $f=\phi(k)$, for some $k \in H^2(\mathbb{D})\ominus \mathcal{N}$.
Now, for any $\phi(h) \in \phi(\mathcal{N})$,
\begin{eqnarray*}
\langle f, \phi(h) \rangle &=& \langle \phi(k), \phi(h) \rangle\\
&=& \frac{1}{4}\bigg(\langle k,h \rangle +  \langle \widehat{k},h \rangle +  \langle k, \widehat{h} \rangle + \langle \widehat{k},\widehat{h} \rangle \bigg)\\
&=& \frac{1}{4}\bigg(\langle k,h \rangle +  \langle \widehat{h},k \rangle +  \langle k, \widehat{h} \rangle + \langle h,k \rangle \bigg) \quad \text{(By Lemma \ref{Lem1})}\\
&= & 0  \quad\text{(since $\widehat{\mathcal{N}}\subset \mathcal{N}$)}
\end{eqnarray*}
which shows that $\phi(H^2(\mathbb{D})\ominus \mathcal{N}) \subset H^2_{\mathbb{R}}(\mathbb{D})
\ominus \phi(\mathcal{N})$. This completes the proof.
\end{proof} 

\end{document}
